\documentclass[a4paper]{article}

% if you need to pass options to natbib, use, e.g.:
%     \PassOptionsToPackage{numbers, compress}{natbib}
% before loading neurips_2022


% ready for submission
%\usepackage{ofu_xai_2022}

\input{math_commands.tex}

% to compile a preprint version, e.g., for submission to arXiv, add add the
% [preprint] option:
%     \usepackage[preprint]{ofu_xai_2022}


% to compile a camera-ready version, add the [final] option, e.g.:
\usepackage[final, nonatbib]{ofu_xai_2022}

% to avoid loading the natbib package, add option nonatbib:
%    \usepackage[nonatbib]{ofu_xai_2022}


\usepackage[utf8]{inputenc} % allow utf-8 input
\usepackage[T1]{fontenc}    % use 8-bit T1 fonts
\usepackage{hyperref}       % hyperlinks
\usepackage{url}            % simple URL typesetting
\usepackage{booktabs}       % professional-quality tables
\usepackage{amsfonts}       % blackboard math symbols
\usepackage{nicefrac}       % compact symbols for 1/2, etc.
\usepackage{microtype}      % microtypography
\usepackage{xcolor}         % colors
\usepackage{longtable}
\usepackage{dsfont}

\usepackage[numbers]{natbib}
%%%%

\title{Penultimate-Layer Feature-Based OOD Detection: Feature Shaping and Its Limits}




\author{%
  Martin Simons\\
  Otto-Friedrich University of Bamberg\\
  Aufbaustr. 11, 96049 Bamberg, Germany\\
  \texttt{martin.simons@stud.uni-bamberg.de} \\[0.5cm]
  Seminar: xAI-Sem-M \\
  Degree: M.Sc. (CitH) \\
  Matriculation \#: 2188932
}


\begin{document}


\maketitle
\def\va{{\bm{a}}}

\begin{abstract}
  The abstract paragraph should be indented \nicefrac{1}{2}~inch on
  both the left- and right-hand margins. Use 10~point type, with a vertical
  spacing (leading) of 11~points.  The word \textbf{Abstract} must be centered,
  bold, and in point size 12. Two line spaces precede the abstract. The abstract
  must be limited to one paragraph.
\end{abstract}


\section{Introduction}
\label{intro}

% Introduction --- slots in after the motivating paragraphs.
% High-level: defines feature shaping and previews the five methods.
% Notation is deferred to the Methodology section.

Feature shaping refers to a family of post-hoc OOD detection methods that act on the
penultimate-layer feature representation of a frozen, pre-trained classifier. Rather
than reading the network's output directly, these methods insert a fixed
transformation that rewrites the feature vector before it reaches the final linear
layer, so that the resulting OOD score separates in-distribution (ID) from
out-of-distribution (OOD) inputs more sharply. The transformation is applied at test
time only and is designed to leave ID classification essentially unchanged; no
weights are retrained. Methods in this family differ in whether they reshape each
feature coordinate independently (\emph{element-wise}) or transform the feature
vector jointly (\emph{whole-vector}), and in how the reshaping rule is obtained.

ReAct~\citep{react} is the simplest: it caps each feature at an upper threshold,
attenuating the abnormally large activations that OOD inputs tend to trigger.
ASH~\citep{ash} instead operates on the whole vector at once, keeping only a small
fraction of the largest activations per input and discarding the rest before
rescaling what remains; unlike ReAct, it draws on no statistics from training data.
VRA~\citep{vra} returns to per-coordinate reshaping but motivates its rule from a
variational argument, suppressing both unusually small and unusually large
activations while amplifying mid-range ones.

Optimal-FS~\citep{optimalfs} approaches the same goal from an optimization
standpoint. Rather than proposing another hand-designed rule, it asks which feature
values contribute most to distinguishing ID from OOD inputs, and derives the
reshaping rule that best exploits them---obtainable in closed form from ID data
alone. This optimization view doubles as a unifying account of the element-wise
methods that came before it: ReAct and VRA can be recovered as approximations to its
optimal element-wise shaping function. The account is, however, specific to
element-wise shaping; whole-vector methods such as ASH lie outside the concrete
problem it solves---a gap its authors leave open, and one this report revisits when
asking what feature shaping is really exploiting.

The fifth method, BLOOD~\citep{blood}, is included as a deliberate contrast rather
than another shaping rule. It leaves the feature vector untouched and instead
measures how smoothly the network transforms its internal representations, on the
premise that a trained model processes ID inputs more smoothly than OOD ones.
Because it reads a property of the representation rather than editing it, BLOOD
isolates the assumption the four shaping methods share---that the OOD signal resides
in the penultimate feature vector itself---and lets the report test whether that
shared premise, rather than any particular reshaping recipe, is what does the work.

\section{Related Work}
\label{related}

\citeauthor{survey} provide a categorization of OOD methods that this paper will follow. According to their taxonomy, feature based OOD detection methods can be differentiated from other approaches by being training agnostic and post-hoc. This categorization contains all approaches that neither need to be implemented in a given model during training nor influence the training process of a given model. Feature based approaches are further categorized as post-hoc. This attribute is assigned to methods that can be applied to pretrained models at inference time. They therefore do not require to update model weights and biases. \citeauthor{survey} differentiate the OOD detection approaches further into feature based, gradient based, output based \dots


% Drops into the existing Related Work section.
% Adjust \subsection -> \section if this is standalone.
\subsection{Positioning within the Survey Taxonomy}
\label{subsec:taxonomy}

We adopt the mechanism-oriented taxonomy of \citet{survey} as the backbone for
situating the methods studied in this report. Its top-level criterion is whether
a method is permitted to intervene in the training process. \emph{Training-driven}
approaches improve OOD separability by modifying the optimization objective or the
training data, whereas \emph{training-agnostic} approaches operate on an
already-trained, frozen model. The cluster examined here lies entirely on the
training-agnostic side: none of the methods alter network weights or the training
procedure.

Within training-agnostic detection, \citet{survey} draw a second distinction
according to whether a method exploits statistical dependencies across the test
stream. \emph{Test-time adaptive} methods update the model or its running
statistics from batches of incoming test data, while \emph{post-hoc} methods score
each input independently, reusing only quantities the frozen network already
computes during a single forward pass. Their lightweight, deployment-friendly
character makes post-hoc detection the dominant practical setting, and it is the
setting to which every method discussed in this report belongs.

Post-hoc detectors are separated further by which model-internal quantity their
scoring function reads: the output/logit vector (e.g.\ MSP, MaxLogit, Energy), a
distance in feature space (e.g.\ Mahalanobis), an input or parameter gradient, the
penultimate feature representation itself, or an estimated density. This report
concentrates on the \emph{feature-based} family, which derives its OOD signal from
the penultimate-layer activation vector $\mathbf{z} \in \mathbb{R}^{M}$ produced
before the linear classifier $\mathbf{W} \in \mathbb{R}^{M \times C}$.

\citet{survey} characterise feature-based methods as differing chiefly in whether
they reshape activations directly, exploit a statistical property of the
representation, or design a new metric over it. The first branch---\emph{feature
shaping}---replaces $\mathbf{z}$ with a transformed $\bar{\mathbf{z}} =
f(\mathbf{z})$, $f : \mathbb{R}^{M} \to \mathbb{R}^{M}$, before the logits are
recomputed, with $f$ chosen so that the resulting score sharpens the ID/OOD gap
while leaving ID classification largely intact. ReAct~\cite{react},
ASH~\cite{ash}, VRA~\cite{vra}, and OptFS~\cite{opt} form a progression
within this branch, moving from heuristic intervention---activation truncation,
top-$K$ pruning with rescaling, and variationally motivated suppression and
amplification---toward Optimal-FS, which casts shaping as an optimization problem
and shows that the earlier methods implicitly approximate a single optimal,
ID-only shaping function.

BLOOD~\cite{blood} also reads the penultimate representation but does not reshape
it. Instead it measures a statistical property of the network's intermediate
transformation---its smoothness, estimated through the feature Jacobian---and is
therefore grouped by \citet{survey} with metric-design methods rather than with
shaping. In this report it serves as a controlled counterpoint that isolates the
shaping cluster's shared assumption: that the discriminative OOD signal resides in
the penultimate activation vector itself, recoverable by a value-to-value
transformation of $\mathbf{z}$.

% Follows \ref{subsec:taxonomy} within Related Work.
\subsection{A Second Axis: Access Requirements}
\label{subsec:access}

The mechanism taxonomy of Section~\ref{subsec:taxonomy} groups methods by
\emph{which} model-internal quantity they read. \citet{blood} argue that a second,
arguably more practically relevant, axis groups methods by the \emph{resources}
they require, and partition detectors into three categories. \emph{Black-box}
methods need only input--output access and are therefore applicable even to models
served behind a product API; the output-based scores MSP, MaxLogit, and Energy are
the canonical examples. \emph{White-box} methods additionally require the model's
weights and architecture but no data, making them directly applicable to
third-party pre-trained checkpoints. \emph{Open-box} methods carry the heaviest
prerequisites: unrestricted access permitting intervention in training, or use of
the training set or a separate OOD validation set.

This axis is orthogonal to the mechanism taxonomy and cuts across the feature-based
family studied here. Although ReAct~\cite{react}, ASH~\cite{ash}, VRA~\cite{vra},
and Optimal-FS~\cite{opt} share a shaping mechanism, they differ in access
requirements: ASH is purely data-free---its top-$K$ pruning needs only the weights,
placing it on the white-box side---whereas ReAct's percentile threshold, VRA's
calibrated cut-offs, and Optimal-FS's closed-form shaping function are all fit to ID
data, placing them in the open-box category. BLOOD~\cite{blood} reads a different
quantity (transformation smoothness via the feature Jacobian) but, requiring only
the weights, is likewise white-box. The clean contrast on this axis is therefore not
BLOOD against the shaping cluster, but the data-free methods (ASH and BLOOD) against
the ID-data-dependent ones (ReAct, VRA, Optimal-FS).

The open-box category also supplies the strongest reference point. Mahalanobis
distance~\cite{mahalanobis}, which fits class-conditional Gaussians to ID training
features and is thus open-box, serves throughout this report as an honest ceiling
rather than a member of either shaping or smoothness paradigm. In the white-box
comparisons of \citet{blood} it outperforms BLOOD in nearly every transformer
setting, and across benchmarks it frequently exceeds both the shaping cluster and
BLOOD---so neither feature-based paradigm is the overall top performer. This
dominance is, however, setting-specific: Mahalanobis is strong on
transformer/text backbones but markedly weaker on ImageNet-scale convolutional
networks, where its class-conditional Gaussian assumption.

\section{Methodology}
\label{sec:methodology}

We consider a classifier pre-trained on in-distribution (ID) data and held frozen.
For an input $\mathbf{x}$, the backbone produces a penultimate feature vector
$\mathbf{z} \in \mathbb{R}^{M}$, which a linear head $\mathbf{W}$ maps to $C$ class
logits. A post-hoc detector reduces the network's response to a scalar score
$s(\mathbf{x}) \in \mathbb{R}$ and flags the sample once that score crosses a
threshold. Detection is thus a thresholding rule
\begin{equation}
G_\lambda : \mathbb{R} \to \{\text{ID}, \text{OOD}\},
\end{equation}
parameterized by a threshold $\lambda$, that labels $\mathbf{x}$ by comparing
$s(\mathbf{x})$ to $\lambda$. The rule $G_\lambda$ is common to every method; the
methods differ only in how the score $s$ is computed from the frozen network and in
its orientation (which side of $\lambda$ denotes ID). Detection quality is reported
with the two standard metrics, AUROC (threshold-free separability) and FPR95 (false
positive rate at $95\%$ true positive rate).

\subsection{Feature Shaping}
\label{subsec:shaping}

Feature shaping inserts a transformation $f : \mathbb{R}^{M} \to \mathbb{R}^{M}$ at
the penultimate layer, reshaping $\mathbf{z}$ before the logits are recomputed and
scored. The detector then factors into a shaping stage and a scoring stage,
\begin{equation}
f : \mathbb{R}^{M} \to \mathbb{R}^{M}, \qquad
s : \mathbb{R}^{C} \to \mathbb{R}, \qquad
s_{\mathrm{shaped}}(\mathbf{x}) = s\big(\mathbf{W} f(\mathbf{z})\big),
\end{equation}
and the sample is flagged by $G_\lambda\big(s_{\mathrm{shaped}}(\mathbf{x})\big)$.
The score $s$ is a standard logit-based score (MSP, MaxLogit, or energy), for which
a lower value indicates OOD; crucially, the choice of $s$ is orthogonal to the
shaping function $f$, a separation Optimal-FS later exploits. Setting
$f = \mathrm{id}$ recovers the unshaped baseline. Shaping functions are
\emph{element-wise} when $f$ acts coordinatewise, $\bar{z}_i = f(z_i)$, and
\emph{whole-vector} when the transformation of one coordinate depends on the others;
this distinction, carried entirely by the single symbol $f$, separates the methods
below.

\subsection{ReAct and VRA}
\label{subsec:react-vra}

ReAct~\citep{react} and VRA~\citep{vra} are both element-wise. ReAct clips each
activation at an upper threshold $c$,
\begin{equation}
f_{\mathrm{ReAct}}(z_i) = \min(z_i, c),
\end{equation}
with $c$ set to a high percentile (typically the $90$th) of the activations
estimated on ID data. The motivation is that OOD inputs tend to produce a few
abnormally large activations; capping them suppresses the OOD score more than the ID
score while leaving most ID activations, which fall below $c$, untouched.

VRA starts from a variational analysis of which operation maximally separates ID and
OOD activations, concluding that the ideal rule suppresses both abnormally low and
abnormally high activations. Its canonical piecewise realization, VRA-P, uses two
ID-estimated quantile thresholds $\alpha < \beta$,
\begin{equation}
f_{\mathrm{VRA\text{-}P}}(z_i) =
\begin{cases}
0 & z_i < \alpha, \\
z_i & \alpha \leq z_i \leq \beta, \\
\beta & z_i > \beta .
\end{cases}
\end{equation}
This strictly generalizes ReAct, recovered as the special case $\alpha = 0$: VRA-P
adds a lower cut that zeroes the abnormally small activations ReAct leaves in place.
Both methods fit their thresholds to ID data and are therefore data-dependent
(cf.\ Section~\ref{subsec:access}). The amplification of mid-range activations
(VRA+) and a feature--logit combined score (VRA++) are further refinements that
depart from pure shaping; we summarize them in the appendix.

\subsection{ASH}
\label{subsec:ash}

ASH~\citep{ash} operates on the whole vector. For each sample it computes the $p$th
percentile $t$ of that sample's own activations, prunes everything below $t$ to
zero, and then treats the surviving activations in one of three ways: ASH-P leaves
them unchanged, ASH-B sets them all to a common constant, and ASH-S---the version we
take as canonical---scales the survivors by a factor that restores the pre-pruning
activation mass,
\begin{equation}
f_{\mathrm{ASH\text{-}S}}(z_i) =
\begin{cases}
z_i \cdot \exp\!\left( s_1 / s_2 \right) & z_i \geq t,\\[2pt]
0 & z_i < t,
\end{cases}
\qquad
s_1 = \sum_j z_j, \quad s_2 = \sum_{j : z_j \geq t} z_j ,
\end{equation}
where $s_1$ and $s_2$ are the sums of activations before and after pruning. The
premise is that the representation is heavily redundant, so most of it can be
discarded and the sparse surviving support---rescaled to compensate---still
classifies ID inputs while sharpening OOD separation. Because the threshold $t$ is
computed per sample from the whole vector, ASH is whole-vector and not reducible to
an element-wise map, and it uses no training statistics, making it data-free
(cf.\ Section~\ref{subsec:access}). An extreme ablation, ASH-Rand, replaces the
surviving values with random ones; we return to it in the discussion, as it bears on
what ASH actually exploits.

\subsection{Optimal-FS}
\label{subsec:optfs}

Optimal-FS~\citep{optimalfs} recasts element-wise shaping as an optimization
problem. Writing the per-coordinate reshaping as a multiplicative factor
$\theta(z_i) = f(z_i)/z_i$, the shaped maximum logit is
\begin{equation}
\bar{\ell}_{\max}(\bar{\mathbf{z}}) = \sum_{i=1}^{M} \theta(z_i)\, w^{\max}_i\, z_i ,
\end{equation}
so the question becomes which feature \emph{value ranges} should be up- or
down-weighted. Optimal-FS partitions the range of feature values into $K$ intervals
$\{[a_k, b_k)\}_{k=1}^{K}$, treats $\theta$ as a piecewise constant
$\hat{\theta}_k$ over each, and collects the per-interval contributions into a
$K$-dimensional interval-specific feature impact (ISFI) vector $\mathbf{I}(\mathbf{z})$,
giving $\bar{\ell}_{\max} = \boldsymbol{\theta}^{\top} \mathbf{I}(\mathbf{z})$. The
shaping vector is then optimized to separate the ID and OOD maximum logits, under a
norm constraint that prevents trivial rescaling,
\begin{equation}
\max_{\boldsymbol{\theta}} \;\;
\boldsymbol{\theta}^{\top}
\Big( \mathbb{E}_{\mathrm{ID}}\big[\mathbf{I}(\mathbf{z})\big]
- \mathbb{E}_{\mathrm{OOD}}\big[\mathbf{I}(\mathbf{z})\big] \Big),
\qquad \text{s.t.} \;\; \|\boldsymbol{\theta}\| = \sqrt{K}.
\end{equation}
Solved with OOD data, this $\boldsymbol{\theta}$ closely matches the element-wise
rules above, recovering them as approximations to a common optimum. Arguing that the
OOD term contributes little beyond the ID direction, Optimal-FS drops it and obtains
a closed-form solution from ID data alone,
\begin{equation}
\boldsymbol{\theta}^{\star} = \frac{\sqrt{K}}
{\big\| \mathbb{E}_{\mathrm{ID}}[\mathbf{I}(\mathbf{z})] \big\|}\,
\mathbb{E}_{\mathrm{ID}}\big[\mathbf{I}(\mathbf{z})\big] .
\end{equation}
Nearly all of the benefit comes from $K \leq 10$, indicating the optimal shaping
function has a simple form (the full ISFI derivation is in the appendix). Two
caveats matter for this report. First, the optimization is specific to element-wise
shaping, so it unifies ReAct, VRA-P, and BFAct but explicitly \emph{does not} capture
whole-vector methods such as ASH, which the authors name as an open problem.
Second, like ReAct and VRA, it depends on ID data (cf.\ Section~\ref{subsec:access}).

\subsection{BLOOD}
\label{subsec:blood}

BLOOD~\citep{blood} reads a different signal from the same region of the network.
Rather than reshaping $\mathbf{z}$, it measures how smoothly the network transforms
its representations, on the premise that a trained model maps ID inputs through
smoother transformations than OOD ones. Smoothness at layer $l$ is the squared
Frobenius norm of that layer's Jacobian $\mathbf{J}_l$,
\begin{equation}
\phi_l(\mathbf{x}) = \big\| \mathbf{J}_l(\mathbf{h}_l) \big\|_F^{2},
\end{equation}
estimated without forming the Jacobian via $N$ pairs of isotropic random probes
$\mathbf{v}_i, \mathbf{w}_i$ (zero mean, unit variance),
\begin{equation}
\hat{\phi}_l(\mathbf{x}) = \frac{1}{N} \sum_{i=1}^{N}
\big( \mathbf{w}_i^{\top} \mathbf{J}_l(\mathbf{h}_l)\, \mathbf{v}_i \big)^{2},
\qquad
\mathbb{E}\big[\hat{\phi}_l(\mathbf{x})\big] = \big\| \mathbf{J}_l \big\|_F^{2},
\end{equation}
the isotropy of the probes being what makes the estimator unbiased. BLOOD uses this
quantity directly as its score, $s(\mathbf{x}) = \hat{\phi}_{L-1}(\mathbf{x})$, with
the \emph{inverted} orientation relative to the shaping cluster: a higher value
indicates OOD. The detector $G_\lambda$ still applies, only with the sides of
$\lambda$ reversed. Note that this score is not expressible in the
shaping composition $s \circ \mathbf{W} \circ f$ of Section~\ref{subsec:shaping}: it
reads the Jacobian and never factors through the head or a shaping map, which is
precisely why BLOOD sits outside feature shaping. The method has two variants,
BLOOD$_M$ averaging $\hat{\phi}_l$ over all layers and BLOOD$_L$ using only the final
transformation, the Jacobian of the head with respect to the penultimate
representation $\mathbf{z}$. The authors focus on BLOOD$_L$, which outperforms
BLOOD$_M$ more often, attributing this to the final layer undergoing the most
training-driven change; this coincides with our penultimate-layer scope, so we take
BLOOD$_L$ as canonical. BLOOD requires only the model's weights and no training data
(cf.\ Section~\ref{subsec:access}).



\section{Experiments}
\label{sec:experiments}

The five methods were never evaluated under a shared protocol, and no honest unified
leaderboard exists: they differ in modality (vision vs.\ text), backbone, OOD-dataset
suite, and scoring function. Rather than force a comparison the literature does not
support, we read the evidence as two internally consistent \emph{islands}---vision
feature shaping and text between-layer smoothness---and treat their incomparability
as itself a finding: these methods target different deployment settings, and their
relative merits are only defined within a setting. All figures below are reproduced
from the cited papers, not independently re-run, and are comparable only within an
island.

\subsection{Vision: the shaping cluster under one protocol}
\label{subsec:exp-vision}

Optimal-FS~\citep{optimalfs} provides the only head-to-head evaluation of the shaping
cluster, re-running ReAct, ASH (P/B/S), VRA-P, and its own method across eight
ImageNet-1k backbones under one protocol. On ResNet-50---the architecture all of these
methods were designed on---the shaping methods are strong and close together (AUROC
$\approx 90.3$ for ASH-S, $88.6$ for VRA-P, $86.5$ for ReAct, $88.6$ for Optimal-FS).
The picture inverts off ResNet-50. On transformer and MLP backbones the most aggressive
shaping rules collapse toward chance: ASH-S falls to $18.5$ AUROC on ViT-B-16 and
$33.9$ on MLP-B, and VRA-P behaves similarly, dragging their eight-backbone averages to
$\approx 38$. ReAct degrades far more gracefully (averaging $81.3$, never falling below
the high $70$s), and Optimal-FS holds across the board ($83.0$ average, $82.7$ on
ViT-B-16). So the cluster's apparent unity is conditional on its home architecture; the
clip-based ReAct and the ID-optimized Optimal-FS generalize, while the percentile-prune
methods do not. On the smaller-scale CIFAR benchmarks, feature shaping as a whole
underperforms simple logit baselines, with Optimal-FS competitive but not dominant. Two
caveats apply: these numbers come from a paper with a stake in the comparison, and the
eight-dataset ImageNet averaging differs from the four-dataset suites used in the
original ASH and ReAct papers.

\subsection{Text: between-layer smoothness}
\label{subsec:exp-text}

BLOOD~\citep{blood} evaluates a disjoint setting---text classification with RoBERTa and
ELECTRA over eight ID datasets---so it neither confirms nor contradicts the vision
results. Among methods requiring no training data (white/black-box), BLOOD$_L$ is the
strongest on most datasets, and more consistently than the all-layer variant BLOOD$_M$;
the authors accordingly focus on BLOOD$_L$, attributing its edge to the final layer
absorbing the most training-driven change. The advantage is clearest on ELECTRA, where
BLOOD$_L$ leads the white/black-box group on all eight datasets. Notably, the two
shaping methods that appear here, ASH and ReAct, are only middling in this setting,
indicating that vision-tuned shaping does not transfer to text. BLOOD's own grouping
also mirrors our access axis (Section~\ref{subsec:access}): ASH is filed as
white/black-box (data-free) and ReAct as open-box (it needs an ID-data percentile).

\subsection{The ceiling, and what the islands show}
\label{subsec:exp-ceiling}

A distance-based open-box method, Mahalanobis distance~\citep{mahalanobis}, bounds both
islands from above. In BLOOD's comparison it outperforms BLOOD in every setup but one
(ELECTRA on TREC), and it is consistently the top scorer on the transformer text
backbones. Its dominance is, however, setting-specific rather than universal: on
ImageNet-scale ConvNets it is weak (AUROC $\approx 55$ on ResNet-50 in VRA's
table~\citep{vra}), where its class-conditional Gaussian assumption is harder to
satisfy. The two islands therefore deliver a consistent structural message rather than
a numerical ranking: feature shaping wins only on its home architecture and modality,
between-layer smoothness occupies a different modality entirely, and a simple
distance-based detector outperforms both wherever its distributional assumptions
hold---so neither shaping nor smoothness is the better paradigm in any setting-free
sense.

\section{Discussion}
\label{sec:discussion}

The five methods share the penultimate layer of a frozen network but were developed
separately, on different architectures and modalities, and never co-evaluated under one
protocol. Placing them together raises three questions: why the same methods succeed on
some architectures and fail on others (Section~\ref{subsec:disc-arch}); how far the
shaping cluster can be formally unified (Section~\ref{subsec:disc-unify}); and whether
the OOD signal lies in the feature values or in the transformation applied to them
(Section~\ref{subsec:disc-carrier}).

\subsection{Architecture dependence}
\label{subsec:disc-arch}

Shaping rules designed on ResNet-50 do not transfer: ASH-S falls from $90.3$ to $18.5$
AUROC from ResNet-50 to ViT-B-16~\citep{optimalfs}. Optimal-FS reports this as a
pattern---for shaping methods, convolutional-network performance is negatively
correlated with performance on other architectures~\citep{optimalfs}. The methods'
mechanisms account for it. ASH's pruning assumes a redundant, sparse representation and
is framed as removing excess features of an over-parameterized
network~\citep{ash}, a ReLU convolutional premise. BLOOD instead reads transformation
smoothness in the upper layers of a transformer, where fine-tuning concentrates and
where ID and OOD representations diverge~\citep{blood}. We read these together as one
point: each method encodes the architecture it was designed on, which is why neither
transfers.

ReAct gives a partial cause. It attributes the abnormal OOD activations to
batch-normalization statistics estimated on ID data and applied to OOD inputs, making
batch normalization the exploited signal rather than an obstacle~\citep{react}. ReAct
also helps under weight and group normalization~\citep{react}, so the effect is not
specific to batch normalization. The architecture dependence is therefore real but its
root cause is unresolved across the cluster.

\subsection{How far the shaping cluster unifies}
\label{subsec:disc-unify}

Optimal-FS defines shaping as optimization over functions
$f : \mathbb{R}^{M} \to \mathbb{R}^{M}$ and proves that ReAct and VRA-P, rewritten via
$\theta(z) = f(z)/z$, approximate the optimal element-wise function~\citep{optimalfs}.
This optimum is simple: most of the benefit occurs at $K \leq 1$ intervals, which prunes
only the tails of the feature distribution~\citep{optimalfs}. This explains why ReAct, a
single clip, stays competitive with the optimization-derived method on the architectures
both were tuned for: little structure beyond the tails remains to recover.

The unification is bounded, and Optimal-FS draws the boundary. Its tractable instance
restricts to element-wise $f$ and, in its own statement, fails to capture whole-vector
methods such as ASH, which it still places inside the general
$\mathbb{R}^{M} \to \mathbb{R}^{M}$ class~\citep{optimalfs}. ASH is thus excluded by the
reduction, not the definition, and the whole-vector extension is named as open future
work~\citep{optimalfs}. Extending this boundary to a case Optimal-FS does not treat, we
observe that BLOOD lies outside the class: its score is the Frobenius norm of a
Jacobian~\citep{blood}, not a map returning a reshaped feature vector, so no extension of
the optimization can absorb it. This gives three tiers: ReAct and VRA-P inside the solved
problem, ASH inside the general class but outside the solved instance, and BLOOD outside
the class. The first gap is one of tractability; the second is categorical.

Whether the whole-vector extension is reachable is open. Two findings bear on it, and
connecting them is our reading rather than a stated result. The element-wise optimum is
coarse~\citep{optimalfs}, and ASH's signal is also coarse: ASH-Rand, which replaces the
surviving activations with random values, still beats the logit baselines, so the signal
lies in the support and the global scale rather than the value profile~\citep{ash}. We
suggest both point the same way---a whole-vector optimum would also be simple, consistent
with Optimal-FS's conjecture~\citep{optimalfs}. The open obstacle is that such a
formulation must optimize which coordinates survive, a combinatorial selection unlike the
element-wise interval weights. The same finding qualifies the ``prior methods approximate
the optimum'' reading for ASH: it is a curve-to-curve match for the element-wise
methods~\citep{optimalfs}, but for ASH the correspondence is to the support-and-scale
operation, not a value curve.

\subsection{Where the signal lives}
\label{subsec:disc-carrier}

Because BLOOD is not a shaping function, it serves as a counterpoint that differs from
the cluster on one variable: what is measured at the top of the network. The four shaping
methods edit the activation vector; BLOOD leaves it unchanged and measures the smoothness
of the transformation applied to it~\citep{blood}. This poses a question neither paradigm
sets out to answer: is the OOD signal a property of the feature values or of the map
applied to them?

We do not resolve which is fundamental. The juxtaposition does narrow the cluster's
premise. BLOOD shows the activation vector is not the only carrier of OOD signal at the
penultimate layer, and ASH-Rand shows that within the cluster the fine values are not the
operative part. The shared assumption---that the signal can be recovered by reshaping the
penultimate activation vector---is therefore narrower than its framing, though not
incorrect.

\section{Conclusion}
\label{conc}


\bibliography{ood_base}
\bibliographystyle{abbrvnat}


%%%%%%%%%%%%%%%%%%%%%%%%%%%%%%%%%%%%%%%%%%%%%%%%%%%%%%%%%%%%%%%%%%%%%%%%%%%%%%%%%%%%%%%%%%%%%%%%%%%%
%% Use of Large Language Models
%%%%%%%%%%%%%%%%%%%%%%%%%%%%%%%%%%%%%%%%%%%%%%%%%%%%%%%%%%%%%%%%%%%%%%%%%%%%%%%%%%%%%%%%%%%%%%%%%%%%
\section*{Use of Large Language Models}


%%%%%%%%%%%%%%%%%%%%%%%%%%%%%%%%%%%%%%%%%%%%%%%%%%%%%%%%%%%%%%%%%%%%%%%%%%%%%%%%%%%%%%%%%%%%%%%%%%%%
%% Declaration of Authorship
%%%%%%%%%%%%%%%%%%%%%%%%%%%%%%%%%%%%%%%%%%%%%%%%%%%%%%%%%%%%%%%%%%%%%%%%%%%%%%%%%%%%%%%%%%%%%%%%%%%%




{\parindent 0cm
%%%%%%%%%%%%%%%%%%%%%%%%%%German%%%%%%%%%%%%%%%%%%%%%%%%%%%%%%
\section*{Declaration of Authorship}
Ich erkläre hiermit gemäß § 9 Abs. 12 APO, dass ich die vorstehende Seminararbeit selbstständig verfasst und keine anderen als die angegebenen Quellen und Hilfsmittel benutzt habe. Des Weiteren erkläre ich, dass die digitale Fassung der gedruckten Ausfertigung der Seminararbeit ausnahmslos in Inhalt und Wortlaut entspricht und zur Kenntnis genommen wurde, dass diese digitale Fassung einer durch Software unterstützten, anonymisierten Prüfung auf Plagiate unterzogen werden kann.\\
\vspace{2\baselineskip}
  
Bamberg, \today

\rule[0.5em]{14em}{0.5pt} \hspace{0.25\linewidth}\rule[0.5em]{14em}{0.5pt}
\vspace{1em}
\hspace{4em} (Ort, Datum) \hspace{0.51\linewidth} (Unterschrift)
}



%%%%%%%%%%%%%%%%%%%%%%%%%%%%%%%%%%%%%%%%%%%%%%%%%%%%%%%%%%%%


%%%%%%%%%%%%%%%%%%%%%%%%%%%%%%%%%%%%%%%%%%%%%%%%%%%%%%%%%%%%


\appendix


\section{Appendix}


% Requires: \usepackage{booktabs} and \usepackage{longtable}
% Replace the placeholder citation keys (react2021, vra2023, ash2023,
% optfs2024, blood2024, goodfellow2016) with your actual .bib keys.

\begin{longtable}{@{}ll@{}}
\caption{Unified notation used throughout this report. \\ 
Symbols are grouped by the work that introduces the corresponding quantity in the OOD-detection
context. Typographic conventions---bold vectors and matrices,
\(\mathds{1}\{\cdot\}\), and \(\top\) for transpose---follow
\citet{goodfellow2016}. Where a symbol is renamed for cross-paper
consistency, the original symbol is given in parentheses.}
\label{tab:notation}\\
\toprule
\textbf{Symbol} & \textbf{Meaning}\\
\midrule
\endfirsthead
\toprule
\textbf{Symbol} & \textbf{Meaning}\\
\midrule
\endhead
\midrule
\multicolumn{2}{r}{\textit{continued on next page}}\\
\endfoot
\bottomrule
\endlastfoot

\multicolumn{2}{@{}l}{\textit{General notation}}\\
\(x\)                     & Input sample\\
\(z\)            & Penultimate-layer feature (activation) vector\\
\(z_i\)                   & \(i\)-th element of \(z\)\\
\(\bar{z}\)      & Feature vector after shaping / rectification\\
\(M\)                     & Penultimate feature dimension\\
\(C\)                     & Number of classes\\
\(\ell_i,\ \ell^{\max}\)  & \(i\)-th logit; max logit\\
\(\mathbf{J}\)            & Jacobian matrix\\
\(\lVert\cdot\rVert_F\)   & Frobenius norm\\
\(\mathds{1}(\cdot)\)   & Indicator (1 if condition holds, else 0)\\
\(\mathbb{E}[\cdot]\)     & Expectation\\
\midrule

\multicolumn{2}{@{}l}{\citet{react}}\\
\(c\)                     & Activation clip threshold, \(\mathrm{ReAct}(z;c)=\min(z,c)\)\\
\(S(\cdot)\)              & OOD scoring function (MSP, MaxLogit, Energy, \dots)\\
\(\lambda\)               & OOD decision threshold\\
\(G_\lambda\)             & OOD detector (in/out rule on \(S\))\\
\midrule

\multicolumn{2}{@{}l}{\citet{vra}}\\
\(p_{\text{in}},\ p_{\text{out}}\) & ID / OOD marginal densities over activations\\
\(g(\cdot),\ g^{*}(\cdot)\)        & Rectification function; variational optimum\\
\(\tau\)                  & Fidelity--gap trade-off coefficient \,(orig.\ \(\lambda\))\\
\(\alpha,\ \beta\)        & Low / high activation thresholds\\
\(\gamma\)                & VRA+ amplification offset on the intermediate band\\
\(\eta_\alpha,\ \eta_\beta\) & ID quantiles defining \(\alpha,\beta\)\\
\midrule

\multicolumn{2}{@{}l}{\citet{opt}}\\
\(\mathbf{w}^{\max}\)     & Weight vector of the predicted (max-logit) class\\
\(\theta(\cdot)\)         & Element-wise shaping function, \(\theta(z):=f(z)/z\)\\
\(\boldsymbol{\theta}\)   & Per-bucket shaping vector, \(\boldsymbol{\theta}\in\mathbb{R}^{K}\)\\
\(I_{[a,b)}(\mathbf{z})\) & Interval-specific feature impact (ISFI) of bucket \([a,b)\)\\
\(a_k,\ b_k\)             & Edges of the \(k\)-th bucket\\
\(K\)                     & Number of buckets \,(\(K=100\))\\
\midrule

\multicolumn{2}{@{}l}{\citet{blood}}\\
\(\phi_l(x)\)             & Score: squared Frobenius norm of the layer-\(l\) Jacobian\\
\(h_l\)                   & Layer-\(l\) representation \,(\(h_{L-1}=\mathbf{z}\))\\
\(N\)                     & Number of random probes \,(orig.\ \(M\))\\
\(\mathbf{v}_i,\ \mathbf{w}_i\) & Gaussian probe / projection vectors\\
\(\mathbf{u}_i\)          & Jacobian--vector product, \(\mathbf{u}_i=\mathbf{J}_{L-1}\mathbf{v}_i\)\\
\midrule

\multicolumn{2}{@{}l}{\citet{ash}}\\
\(p\)                     & Pruning percentile\\
\(t\)                     & Per-sample threshold, \(t=\mathrm{Perc}_p(\mathbf{z})\)\\
\(m_i\)                   & Binary keep-mask, \(m_i=\mathds{1}\{z_i\ge t\}\)\\
\(n\)                     & Number of surviving activations\\
\(s_1,\ s_2\)             & Pre- / post-pruning activation sums\\

\end{longtable}

\end{document}
